% ============================================================
% Section 152: 一维盒中电子对器壁的压力
% ============================================================
\section{量子力学：一维盒中电子对器壁的压力}
\label{sec:particle-box-pressure}

\begin{modelbox}[题目：一维无限深势阱中电子对器壁的平均力]
一个电子被禁闭在一维盒子中，并处于基态上。盒宽 $a=10^{-10}\,\text{m}$，电子能量为 $38\,\text{eV}$。计算：

(1) 电子在其第一激发态的能量；

(2) 当电子处于基态时盒壁所受的平均力。

提示：平均力计算利用 Hellmann-Feynman 定理。
\end{modelbox}

\begin{tipbox}[考点快速分析]
\begin{itemize}
    \item \textbf{一维无限深势阱}：$E_n=\dfrac{n^2\pi^2\hbar^2}{2ma^2}$
    \item \textbf{能级比例}：$E_n\propto n^2$，第一激发态 $E_2=4E_1$
    \item \textbf{Hellmann-Feynman 定理}：$\langle F\rangle=-\dfrac{\partial E}{\partial b}$，$b$ 为器壁位置参数
    \item \textbf{经典类比}：$F_{\text{cl}}=2E/a$，与量子结果一致
\end{itemize}
\end{tipbox}

\begin{derivationbox}[标准解题过程]
\textbf{(1) 第一激发态能量}

基态对应 $n=1$，$E_1=38\,\text{eV}$。第一激发态 $n=2$：
\begin{equation}
\boxed{E_2=4E_1=4\times38=152\,\text{eV}.}
\end{equation}

\textbf{(2) 基态时盒壁所受的平均力}

设盒子左边固定于 $x=-a/2$，右边壁位于 $x=b$（初始 $b=a/2$）。盒子宽度为 $b+a/2$。基态能量：
\[
E_1(b)=\frac{\pi^2\hbar^2}{2m(b+a/2)^2}.
\]

由 Hellmann-Feynman 定理，粒子对右边壁的平均力：
\[
\langle F\rangle=-\frac{\partial E_1}{\partial b}\bigg|_{b=a/2}=\frac{\pi^2\hbar^2}{ma^3}=\frac{2E_1}{a}.
\]

代入数值：
\begin{equation}
\boxed{\langle F\rangle=\frac{2\times38\,\text{eV}}{10^{-10}\,\text{m}}=7.6\times10^9\,\text{eV/cm}\approx1.2\times10^{-7}\,\text{N}.}
\end{equation}

\textbf{经典类比}：经典粒子以速度 $v$ 在盒内往返，周期 $T=2a/v$。每次撞击右壁动量改变 $2mv$，平均力 $F_{\text{cl}}=2mv/T=mv^2/a=2E/a$，与量子结果完全一致。
\end{derivationbox}

\begin{formulabox}[答案汇总]
\begin{center}
\begin{tabular}{ll}
\toprule
\textbf{物理量} & \textbf{数值} \\
\midrule
(1) 第一激发态能量 $E_2$ & $152\,\text{eV}$ \\
(2) 盒壁平均力 $\langle F\rangle$ & $2E_1/a\approx1.2\times10^{-7}\,\text{N}$ \\
\bottomrule
\end{tabular}
\end{center}
\end{formulabox}

\begin{insightbox}[物理图像]
\begin{itemize}
    \item Hellmann-Feynman 定理将力学量的期望值与能量对参数的导数联系起来，是计算量子系统中广义力的有力工具。
    \item 量子结果与经典图像的一致性体现了对应原理：当量子数很大时，量子力学的预言回归经典结果。
    \item 在定态下 $\langle\partial V/\partial x\rangle=0$，说明左壁与右壁受力大小相等、方向相反，否则盒子会自行加速。
\end{itemize}
\end{insightbox}
